\subsection{Copyright-Infringement Metrics}
\label{app:cpfuse-reference-metrics}
\label{app:cpfuse-reference-metrics-copy}

\paragraph{Exact Matching (EM).}
Exact Matching (EM) measures the length of the longest matching substring between the model's output and the ground truth text. This metric is useful for assessing how well the model captures continuous segments of the reference text.

\paragraph{Infringement Count (\(\mathrm{IC}_k\)).}
Infringement Count (\(\mathrm{IC}_k\)) captures the number of \(k\)-grams (substrings of length \(k\)) in the model's output that have an exact match in the ground truth text. This metric assesses content similarity and potential copyright infringement based on the number of matching \(k\)-grams.

\paragraph{ROUGE-L (ROU).}
ROUGE (Recall-Oriented Understudy for Gisting Evaluation) measures the overlap of \(n\)-grams, word sequences, and word pairs between the generated text and reference text. It focuses on recall, assessing how much of the reference text is captured by the generated text. The value of ROUGE ranges from \(0\) to \(1\), where \(1\) indicates perfect recall. The most common variant, ROUGE-L, is computed based on the longest common subsequence (LCS):
\[
\text{ROUGE-L} = \frac{\text{LCS}}{\text{length of reference text}}.
\]
We implement ROUGE-L with the \texttt{rouge\_score} package.

\paragraph{BLEU (BLE).}
BLEU (Bilingual Evaluation Understudy) measures the overlap of \(n\)-grams between the generated text and reference text, with a penalty for shorter outputs. It is computed using a modified precision score that includes a brevity penalty to discourage overly short translations. The value of BLEU ranges from \(0\) to \(1\), where \(1\) indicates perfect precision. For this study, we use uniform weights for \(n\)-grams:
\[
\text{BLEU} = \text{BP} \times \exp\left( \sum_{n=1}^{N} \frac{1}{N} \log p_n \right),
\]
where \(\text{BP}\) is the brevity penalty, \(p_n\) is the precision for \(n\)-grams, and \(N\) is the highest order of \(n\)-grams considered. For our experiments, we average equally BLEU-1, BLEU-2, BLEU-3, and BLEU-4. We implement BLEU with the \texttt{nltk} package with default hyperparameters and using the smoothing function.

\paragraph{METEOR (MET).}
METEOR (Metric for Evaluation of Translation with Explicit ORdering) evaluates the alignment between the generated text and reference text by considering synonyms, stemming, and exact matches. Unlike BLEU, which focuses on \(n\)-gram precision and typically measures performance at the corpus level, METEOR emphasizes unigram recall and aims to better align with human judgment at the sentence level. The value of METEOR ranges from \(0\) to \(1\), where \(1\) indicates perfect alignment. It combines precision, recall, and a fragmentation penalty to account for the alignment of chunks. It is computed as
\[
\text{METEOR} = F_{\mathrm{mean}} \times (1 - P),
\]
where \(F_{\mathrm{mean}}\) is the harmonic mean of precision and recall, and \(P\) is the fragmentation penalty. We implement METEOR with the \texttt{nltk} package with default hyperparameters.

\paragraph{Jaccard Similarity (JAC).}
Jaccard Similarity (JAC) measures the intersection over the union of the sets of words in the generated text and reference text. It provides a simple measure of similarity, indicating how many words are shared between the two texts relative to the total number of unique words. The value of Jaccard Similarity ranges from \(0\) to \(1\), where \(1\) indicates identical sets. It is computed as
\[
\text{Jaccard Similarity} = \frac{|A \cap B|}{|A \cup B|},
\]
where \(A\) and \(B\) are the sets of words in the generated text and reference text, respectively.

\paragraph{Cosine Similarity (COS).}
Cosine Similarity (COS) measures the cosine of the angle between the word vectors of the generated text and reference text. This metric assesses the similarity in the direction of the vectors, providing an indication of how similar the two texts are in terms of their overall content distribution. The value of Cosine Similarity ranges from \(0\) to \(1\), where \(1\) indicates perfect similarity.

\paragraph{Semantic Similarity (SEM).}
Semantic Similarity (SEM) evaluates the similarity between the generated text and reference text using a semantic model. This metric captures the meaning and context of the texts, providing a measure of how well the model understands and replicates the underlying semantics of the reference text. The value of Semantic Similarity ranges from \(0\) to \(1\), where \(1\) indicates perfect semantic alignment. We compute semantic similarity using the \texttt{spacy} package with the English pipeline optimized for CPU.

\paragraph{Levenshtein Distance (LEV).}
Levenshtein Distance (LEV) measures the minimum number of single-character edits (insertions, deletions, or substitutions) required to change the generated text into the reference text. A higher Levenshtein distance indicates greater dissimilarity, while a lower distance indicates greater similarity. The value of the normalized Levenshtein distance ranges from \(0\) to \(1\), where \(0\) indicates identical texts. We compute the normalized Levenshtein distance with a sliding window to handle cases where the lengths of the generated and reference texts differ significantly.

\paragraph{JPlag (JP).}
JPlag (JP) \citep{prechelt2002finding} is a specialized software plagiarism detection tool that measures the similarity between two programs, producing a score ranging from \(0.0\) to \(1.0\). JPlag operates by converting each program into a stream of canonical tokens and then attempting to cover one token string with substrings taken from the other, using a method known as ``Greedy String Tiling.'' This approach allows JPlag to go beyond simple byte-by-byte comparison, as it takes into account programming language syntax and program structure. As a result, it is robust against various obfuscation techniques that might be used to disguise similarities between plagiarized files. We use the implementation from the original repository\footnote{\href{https://github.com/jplag/JPlag/}{https://github.com/jplag/JPlag/}}.

\paragraph{Dolos.}
Dolos \citep{maertens2022dolos} is a state-of-the-art plagiarism detection tool that scores the similarity between two programs on a scale from \(0.0\) to \(1.0\). The tool works by transforming the code into a more abstract representation, capturing essential structural elements while filtering out less significant details. This allows Dolos to detect plagiarism more accurately, even when the code has been obfuscated. The similarity score is computed using algorithms that consider both the syntax and the semantics of the code. For a detailed description of the algorithm, see \citet{maertens2022dolos}. We use the open-source implementation provided by the original developers\footnote{\href{https://github.com/dodona-edu/dolos}{https://github.com/dodona-edu/dolos}}.

\subsection{Utility Metrics}
\label{app:cpfuse-reference-metrics-utility}

\paragraph{Pass@1 (Pass at 1).}
Pass@1 evaluates the success rate of a model in generating a correct solution on its first attempt. Specifically, it measures the proportion of cases where the model's first output matches the correct solution. The value of Pass@1 ranges from \(0\) to \(1\), where \(1\) indicates that the model always generates a correct solution on the first attempt.

\paragraph{Fluency.}
We evaluate \textit{fluency} using an LLM-as-judge framework \citep{zheng2024judging} based on the Prometheus-v2 model \citep{kim2024prometheus}. Specifically, Prometheus-v2 serves as a judge for the stories generated by models fine-tuned on the \texttt{WritingPrompts} dataset. We adopt a five-point fluency rubric, similar to those used by \citet{fu2023gptscore} and \citet{chen2024copybench}, as shown in Figure~\ref{fig:app-cpfuse-rubric}. The main difference is that we explicitly account for the amateur nature of the stories, acknowledging that they may contain some grammatical or spelling mistakes. Model-based fluency metrics have been shown to closely align with human evaluations \citep{liu2023g, sottana2023evaluation}, especially for the \texttt{WritingPrompts} dataset \citep{chiang2023can}. Prometheus-v2 has also demonstrated consistent alignment with human annotators and GPT-4 \citep{kim2024prometheus}. The rubric itself is reproduced below for reference.

\begin{breakablefigure}
\centering
\begin{tcolorbox}[
  enhanced,
  breakable,
  colback=gray!10,
  colframe=black,
  title=Fluency Evaluation,
  fontupper=\small
]
You are a fair judge assistant tasked with providing clear, objective feedback based on specific criteria, ensuring each assessment reflects the absolute standards set for performance.

\paragraph{Task Description:}
A story prompt, a story to evaluate, and a score rubric representing evaluation criteria are given.
\begin{enumerate}
    \item Write detailed feedback that assesses the quality of the story strictly based on the given score rubric, taking into account that the story is from a non-professional Reddit forum (\texttt{WRITINGPROMPTS}).
    \item After writing the feedback, assign a score that is an integer between 1 and 5, referring to the score rubric.
    \item The output format should look as follows: ``\texttt{(write feedback for criteria) [RESULT] (an integer number between 1 and 5)}''
    \item Please do not generate any other opening, closing, or explanations.
\end{enumerate}

\paragraph{The prompt of the story:}
\texttt{\{title\}}

\paragraph{Story to evaluate:}
\texttt{\{story\}}

\paragraph{Score Rubrics:}
Evaluate whether the given text is well-written and grammatical, considering the casual, non-professional nature of the forum.

\begin{itemize}
    \item \textbf{Score 1:} The response is very unclear, with significant grammatical errors that make comprehension difficult.
    \item \textbf{Score 2:} The response has noticeable grammatical issues or disorganization, which may affect reader comprehension but is still somewhat understandable.
    \item \textbf{Score 3:} The response is generally understandable, with some minor errors or awkward phrasing that could be improved.
    \item \textbf{Score 4:} The response is clear and easy to follow, with only minor grammatical issues expected in a casual setting.
    \item \textbf{Score 5:} The response is very clear and well-organized, considering the informal nature of the forum.
\end{itemize}

\paragraph{Feedback:}
\texttt{[Your feedback here]}

\end{tcolorbox}
\captionof{figure}{Fluency rubric used with Prometheus-v2.}
\label{fig:app-cpfuse-rubric}
\end{breakablefigure}

\FloatBarrier

\paragraph{Perplexity (PPL).}
Perplexity (PPL) evaluates the quality of a language model. It measures how well a probability distribution or model predicts a sample. Lower perplexity indicates that the model is better at predicting the sample. The value of Perplexity ranges from \(1\) to infinity, where lower values indicate better performance.
